\section{Linear programming by pricing}\label{sec:lp}

For a subspace \(C\subseteq\RR^N\), consider the two nonnegative cones in
\(C\) and \(C^\perp\). At each coordinate \(i\), exactly one of these cones
contains a vector positive at \(i\). In the homogenized slack space of a
linear program, the two possibilities at the homogenizing coordinate
give a feasible point and a certificate of infeasibility, respectively.
We first determine which cone contains the witness, then recover the
witness itself with inclusion-minimal support.

Call \(C\) the positive side and \(C^\perp\) the negative side.
Orthogonality excludes simultaneous witnesses at \(i\). If the system
\(v\in C,\ v\ge e_i\) is infeasible, Farkas' lemma gives a nonnegative
vector in \(C^\perp\) positive at \(i\), proving that one side always
has a witness. Moreover, a witness \(u\) of inclusion-minimal support
is a circuit. If there were a second independent vector on its support
in the same subspace, a linear combination would give a nonzero
direction supported there and vanishing at \(i\). Move from \(u\)
along that direction, choosing the sign that decreases some entry,
until the first entry vanishes. This yields a witness with smaller
support. For the relation with elementary vectors and the classical
theorems of Minty and Tucker, see~\cite[Theorem~6 and its corollary]{Rock}
and~\cite{Tuc}.

\subsection{Reciprocal scaling and visibility}

To use Theorem~\ref{thm:prices}, we need a label at each coordinate of
each queried price vector. Scale coordinate \(i\) on the positive side
by a positive factor and on the negative side by its reciprocal. This
preserves orthogonality and can bring a small circuit entry into view.
The following theorem chooses the effective scaling for each query
and describes its effect using a rational matrix of short bit length.

Write \(\Slots=[N]\times\{-1,+1\}\) and abbreviate the slots
\((i,\pm1)\) by \(i\pm\). A \emph{reciprocal scaling} is a positive
diagonal matrix \(D\) on these slots with \(D_{i+}D_{i-}=1\) for all
\(i\). For row vectors, scaling acts by right multiplication. The
scale for a query \(p\) will make every slot maximizing surplus visible,
even when the factors \(2^{-p_i},2^{p_i}\) would be too costly to represent.

\begin{theorem}[Calibration]\label{thm:cal}
Given an ordered row basis \(B_0\) of a real subspace \(W\)
on the paired slots of \([N]\), \(N\ge1\), and an integer precision
parameter \(P\ge1\), there is a deterministic procedure answering
online price queries \(p\in\QQ^N\) by rational matrices \(B\).
On a query it may instead declare a restart.

At any time, all answers since the most recent restart, or since
initialization, admit the same assignment of demands. For each circuit
\(u\) of \(W\) the assignment gives an integer \(a_s\in[1,R]\) at each
slot \(s\) in its support.
For each such answer \(B\) at prices \(p\),
there is a reciprocal scaling \(D\) with the following properties.
\begin{itemize}
 \item \(B\) approximates a true row basis \(B_*=L B_0D\) of \(WD\)
 to within \(2^{-P}\) entrywise, with \(L\) lower triangular with
 positive diagonal. Entries of \(B_*\) have absolute value at most
 \(2\). Any vector in \(WD\) of sup norm one has row coefficients
 in this basis of absolute value at most \(2^{100N^2}\).
 \item For every circuit \(u\), at each slot \(s=(i,h)\) maximizing
 surplus \(a_s-hp_i\) on the support of \(u\),
 \[
       |(uD)_s|\ \ge\ \tfrac13\|uD\|_\infty .
 \]
\end{itemize}

The integer bound \(R=2^{\poly(N,P)}\) is available at initialization,
and the total number of restarts is polynomial in \(N,P\). The
procedure is strongly polynomial when the value of \(P\), the number
of queries, and the bit lengths of the prices count as additional
discrete input size. It returns the answers in explicit binary form.
\end{theorem}

The algorithm receives the approximate basis \(B\). The true basis
\(B_*\), the scaling \(D\), and the demands certify the properties of
that answer. The triangular factor \(L\) will also preserve the row
order and orientation needed in Section~\ref{sec:quadratic}.

\subsection{From approximate labels to exact witnesses}

Embed the two sides in \(W=C\times C^\perp\), with \(C\) on the positive
slots and \(C^\perp\) on the negative slots. A witness circuit of either
side is still a circuit in \(W\) after extension by zero on the other
side. The quadratic form
\[
 Q(y)=\sum_{i=1}^N y_{i+}y_{i-}
\]
vanishes on \(W\) and on every reciprocally scaled space \(WD\).
This identity lets us label a coordinate even though calibration returns
only an approximate basis.

\paragraph{Testing one side.}
Apply Theorem~\ref{thm:cal} to a row basis of \(W\), taking
\begin{equation}\label{eq:labelpars}
 H=2^{100N^2},\qquad \epsilon=H^{-10},\qquad 2^{-P}=H^{-100}.
\end{equation}
Run the algorithm of Theorem~\ref{thm:prices} with the resulting bound
\(R\). At each price query, request the calibrated basis. If calibration
declares a restart, begin a fresh price search, retaining the state of
the calibration procedure. Otherwise let \(B\) be the matrix returned.
Form the convex set
\begin{equation}\label{eq:box}
 \mathcal B(B)=\{y=\xi B:\|\xi\|_\infty\le H,\quad
                     -\epsilon\le y_s\le2\quad(s\in\Slots)\}.
\end{equation}
Give label \(+\) at \(i\) if some \(y\in\mathcal B(B)\) has
\(y_{i+}\ge1/4\), and give label \(-\) otherwise.
Each test is linear feasibility on explicit short rational data, so
ordinary bit-polynomial LP suffices~\cite{Khach,GLS}.

To see why testing the positive slot suffices, write a candidate
as \(y=\xi B\) and compare it with \(y_*=\xi B_*\) in the true
scaled space. There are at most \(2N\) rows, so
\(\|y-y_*\|_\infty\le 2NH2^{-P}<\epsilon^2\).
Since \(Q(y_*)=0\) and the entries of \(y\) belong to
\([-\epsilon,2]\), we have \(Q(y)\le5N\epsilon^2\le H^{-1}\).
The following observation will also apply in quadratic programming.

\begin{lemma}\label{lem:mid}
Let \(K\) be a convex set of vectors on the paired slots, with
\(-\epsilon\le y_s\le2\) and \(Q(y)\le H^{-1}\) for every \(y\in K\).
For each \(i\), at most one of its two slots attains \(1/4\) or more
anywhere in \(K\).
\end{lemma}
\begin{proof}
If a vector attaining the threshold exists at each slot, their
midpoint has both entries at least \((1/4-\epsilon)/2\ge1/10\).
Every other pair contributes at least \(-2\epsilon\). Thus the
midpoint belongs to \(K\) and has
\(Q\ge10^{-2}-2N\epsilon>H^{-1}\), a contradiction.
\end{proof}

\paragraph{Deciding all coordinates.}
In the final run between restarts, calibration supplies a common
assignment of positive demands to every circuit.
Fix a nonnegative circuit \(u\) and consider any query \(p\). If the
slot \(s=(i,h)\) maximizes the surplus of this circuit, its normalized
scaled vector
\[
 y_*=uD/\|uD\|_\infty=\xi B_*
\]
has \((y_*)_s\ge1/3\) and admits coefficients with
\(\|\xi\|_\infty\le H\). All entries of \(y_*\) lie in \([0,1]\).
The same approximation estimate shows that \(\xi B\) belongs to
\(\mathcal B(B)\) and attains at least \(1/4\) at \(s\). For \(h=+1\)
the positive test therefore succeeds. For \(h=-1\) Lemma~\ref{lem:mid}
forces the positive test to fail. In either case the returned label at
\(i\) equals \(h\).

Thus the slots and demands of each nonnegative circuit form a
consistent hypothesis for Theorem~\ref{thm:prices}. The returned
prices \(x\) satisfy \(hx_i\ge a_s\ge1\) at \emph{every} slot of that
circuit. Every coordinate has a witness circuit, so \(\sgn(x_i)\)
identifies the side of every coordinate. Calibration restarts only
polynomially often. In each run Theorem~\ref{thm:prices} bounds
the number of queries and their bit lengths by a polynomial in
\(N,\log(R+1)\). The calibrated answers and the rational label tests
then have strongly polynomial total cost by Theorem~\ref{thm:cal}.

\paragraph{Recovering one circuit.}
Fix a target coordinate \(i\) and take the subspace on its indicated
side. For each coordinate \(j\ne i\) in turn, tentatively impose
\(v_j=0\) on the current subspace. Keep the restriction exactly when
the side decision for this smaller subspace says it still contains
a nonnegative vector positive at \(i\). A rejected restriction can
never become feasible later, since each remaining step shrinks the
subspace. Consequently every final witness is positive on every
unrestricted coordinate. The direction argument at the start of this
section now shows that the final subspace is a line. Linear algebra
recovers its generator, oriented to be positive at \(i\), and hence
a witness circuit in the original subspace. At each trial we compute
a basis from the original input and the retained zero coordinates.
The bit bounds are given in Section~\ref{sec:arithmetic}.

\subsection{Certificates for a linear program}

\begin{proof}[Proof of Theorem~\ref{thm:lp}]
For \(Ax\le b\) use the homogenized slack space and its orthogonal
complement,
\[
 \begin{aligned}
 C      &=\{(\alpha b-Ax,\alpha): x\in\RR^n,\ \alpha\in\RR\},\\
 C^\perp&=\{(y,-b^\top y): A^\top y=0\}.
 \end{aligned}
\]
Recover a witness circuit at the last coordinate. On the positive
side, divide the witness by its last coordinate and solve for \(x\)
by linear algebra. This gives \(Ax\le b\). On the negative side,
the witness gives \(y\ge0\), \(A^\top y=0\), and \(b^\top y<0\).
The codimension of \(C^\perp\) is \(1+\rank A\). A circuit in a
space of codimension \(k\) has at most \(k+1\) entries in its
support, since its support indexes a minimally dependent set of
columns of a defining matrix. The last entry of this circuit is
positive, leaving at most \(1+\rank A\) entries for \(y\).

This also solves feasibility with both strict and weak inequalities.
Homogenize with a variable \(\alpha\), require \(\alpha\ge1\), and
require each slack belonging to a strict inequality to be at least
one. A solution of the original mixed system can be scaled to meet all
these lower bounds. Dividing by \(\alpha\) gives the converse.

For a feasible program, test \(Ar\le0,\ c^\top r>0\) by this
procedure. A solution, together with the feasible point already
found, certifies unboundedness. Otherwise LP duality gives finite
primal and dual optima with a common value \(v\). Strict
complementarity~\cite{GT} ensures that some optimal pair satisfies
\(b-Ax+y>0\). For completeness, take a coordinate \(i\) whose slack
is zero at every primal optimum. Maximize that slack subject to
\(Ax\le b\) and \(c^\top x\ge v\). Duality for this auxiliary
program, which has value zero, gives \(u\ge0\) and \(\lambda\ge0\)
with
\[
 A^\top(u+e_i)=\lambda c,\qquad b^\top(u+e_i)=\lambda v .
\]
If \(\lambda>0\), dividing \(u+e_i\) by \(\lambda\) gives a dual
optimum positive at \(i\). If \(\lambda=0\), adding \(u+e_i\) to
any dual optimum has the same effect. Averaging the primal witnesses
for the coordinates admitting positive slack, and the dual witnesses
for those with necessarily zero slack, gives a strictly complementary
pair. If one of these collections is empty, any optimum on that side
suffices.

The algorithm recovers the pair directly, without knowing \(v\).
Solve the single mixed linear feasibility problem
\[
 Ax\le b,\quad y\ge0,\quad A^\top y=c,\quad
 b^\top y=c^\top x,\quad b-Ax+y>0.
\]
The equality makes the duality gap \((b-Ax)^\top y\) zero. This
returns the optimal pair with the required strict complementarity.
\end{proof}

\subsection{Arithmetic and bit lengths}\label{sec:arithmetic}

The bounds used above apply to one uniform program for all input dimensions.
Our control operations on explicit strings (indices or
short rationals generated by the algorithm)
are charged full bit cost, and can use any polynomial-time
bit algorithm on strings of polynomial length in the dimension.
We take standard explicit input with dimensions and a nonzero end marker
(so size is at least the dimensions, also in the bit bounds).
To simulate the algorithms on the real machine of~\cite{BSS}, store
each explicit bit in a cell. Even sequential memory scans that
increment or reset indices and copy cells incur only polynomial
overhead. All polynomial bounds are fixed and effective. When such a
bound is used as an algorithmic parameter, one can fix it by
composing the estimates in the proof.
Logarithms are to base two unless indicated by \(\ln\).

We can implement all scalar arithmetic using pairs
\((u,v)\), meaning \(u/v\) (\(v>0\)),
and perform only ring operations and comparisons aside from
possible final divisions.
This representation admits two input conventions.
\begin{itemize}
  \item Input a scalar \(x\) directly, starting with \((x,1)\).
    Measure bit lengths on rational data by \emph{reduced} fractions
    (of the entries \(u,v\) themselves, and similarly for intermediates).
  \item Give each rational scalar as numerator and denominator
    (integer scalars). Compute with integer pairs and return the
    resulting unreduced pair for each output scalar.
    Measure bit lengths by the full sizes of these integers.
    This uses only the most restrictive division convention
    in~\cite[\S1.3]{GLS}.
    See also~\cite{DKNOV} for a discussion of arithmetic models.
\end{itemize}
These bounds give a polynomial-time Turing simulation. A simulation
may reduce fractions in the first model and store the integer pairs as given
in the second. Neither algorithm
itself uses fraction reduction in scalar arithmetic. Short rationals \emph{given or constructed in
bits} may be freely placed in scalar computations by setting up their
numerators and denominators bit by bit.

To cover both versions, on rational data take as a pair's log height
the log of the least bound on absolute numerators and denominator
for writing \(u,v\) with one positive denominator.
All results and intermediates must have log height polynomial
in total input length (including bit sizes of any input pairs).
The usual pair operations incur log height at most the sum for their
arguments plus one (take the product of their two common denominators),
including for cross-products when comparing.
This gives the bit bound for arithmetic by polynomial-size formulas
(size counting all leaf occurrences).

We use ordinary linear algebra on explicit real matrices.
We give the determinant routine explicitly because repeated unreduced
fraction arithmetic in elimination could increase the bit lengths
(see also~\cite[\S1.4]{GLS}).
Multiply the given \(k\times k\) matrix by the product \(D\)
of all pair-denominators \(v\), forming each new entry by multiplying its numerator
by all denominators except its own. Thus the resulting matrix is
\(M=D B\), where \(B\) denotes the input matrix. Compute \(k!\det M\)
by Newton's identities. If \(c_i\) is the coefficient of \(t^i\)
in \(\det(I+tM)\), the recurrence is
\[
 i!c_i=\sum_{j=1}^i (-1)^{j-1}
 \frac{(i-1)!}{(i-j)!}
 \big((i-j)!c_{i-j}\big)\tr(M^j),\qquad 1\le i\le k.
\]
Use these division-free ring operations (factorial ratios are
products) and return \((k!c_k,k!D^k)\).
On rational data all entries of \(M\) and \(D\) have
a common denominator \(U\) of polynomial bit length with
polynomial-bit numerators. The degree \(i\) results here
can use denominator \(U^i\), with polynomial-bit numerators
(by determinant expansion for the \(c_i\), and similarly for
matrix powers and then partial results in the recurrence).
The algorithm applies the recurrence to the scalar entries of \(M\)
using ring arithmetic. The common denominator \(U\) serves only to bound
reduced bit lengths in the first input model. In the integer-pair model
all entries of \(M\) are integers, so \(U=1\).

Use bordered minor tests to pivot for rank (an invertible minor of
order less than the matrix rank can be enlarged, by the Schur
complement), and Cramer's rule for solving and obtaining bases.
A fixed number of successive linear-algebra computations and
polynomial-size formulas therefore preserves polynomial bit length.
In witness recovery the number of coordinate restrictions can grow with
the dimension. Each basis there is computed from the original input and
the current list of zero coordinates, which keeps the number of successive
linear-algebra computations on the scalars bounded.

\paragraph{Ordered fields.}
For any fixed input shape, unroll the field algorithm to its uniform
step bound. Each computation branch and its output conditions can
be expressed in the language of ordered rings, with division by a
nonzero denominator expressed by multiplication. Tarski's transfer
principle thus gives correctness with the same step bound over every
real closed field. For an arbitrary ordered field, run the argument in
a real closure. All operations on inputs from the original field stay
in that field, so the same certificates are obtained there.
